\documentclass[11pt]{article}

\usepackage[margin=1in]{geometry}
\usepackage{times}
\usepackage{titlesec}
\usepackage{setspace}
\usepackage{array}

\setlength{\parskip}{0.3em}
\setlength{\parindent}{0pt}

\titlespacing*{\section}
{0pt}{0.6em}{0.3em}

\titlespacing*{\subsection}
{0pt}{0.45em}{0.2em}
\titlespacing*{\subsubsection}
{0pt}{0.35em}{0.15em}

\begin{document}

\begin{center}
    {\Large \textbf{ECE51216: Digital Systems Design Automation}}\\
    {\normalsize Spring 2026}\\[0.4em]
    {\Large \textbf{Final Project Report}}\\[0.3em]
    {\normalsize Option 1: Build}\\[0.6em]
    \textbf{Soomin Lee, Dogyu Ryu}
\end{center}

\section{Motivation and Rationale}

The goal of this project was to build a SAT solver for CNF formulas in DIMACS
format and to study how solver behavior changes when stronger heuristics and
data structures are added on top of a basic DPLL implementation. SAT solving is
a natural topic for this project because it is both algorithmically simple at
the baseline level and highly sensitive to implementation choices. A plain DPLL
solver is easy to understand, but its performance depends heavily on how Boolean
Constraint Propagation (BCP) is implemented and how branch literals are chosen.

For that reason, our rationale was to build the solver in stages. We first
implemented a correct recursive DPLL baseline, then added one improvement that
changes propagation cost and one improvement that changes decision quality:

\begin{itemize}
    \item \textbf{watched literals} for more efficient Boolean Constraint
    Propagation (BCP),
    \item \textbf{DLIS} (Dynamic Largest Individual Sum) for decision literal
    selection.
\end{itemize}

The final system supports three directly comparable variants:

\begin{itemize}
    \item \texttt{baseline}: baseline branch selection + baseline propagation,
    \item \texttt{watched}: baseline branch selection + watched-literal
    propagation,
    \item \texttt{dlis+watched}: DLIS branch selection + watched-literal
    propagation.
\end{itemize}

The project remains a \textbf{DPLL solver}, not a CDCL solver. Therefore,
advanced CDCL-specific techniques such as conflict clause learning,
non-chronological backtracking, clause deletion, or restarts were \textbf{not}
implemented in the final code. The heuristics actually implemented and
evaluated are DPLL recursion, BCP, watched literals, and DLIS. This design
choice kept the implementation focused and made the performance comparison easy
to interpret.

Accordingly, the pseudo-code below describes the basic DPLL algorithm and every
major heuristic implemented in our solver. We do not provide pseudo-code for
conflict clause learning or non-chronological backtracking because those
features were outside the final implementation.

\section{Algorithms and Data Structures Chosen}

\subsection{Baseline DPLL}

The baseline solver follows the standard DPLL structure. It repeatedly performs
unit propagation, checks whether the formula is solved or contradictory, and if
not, chooses a branching literal and recursively explores the two possible
truth assignments.

\begin{quote}
\small
\ttfamily
\raggedright
DPLL(F, A, heuristic, propagation)\\
\hspace*{1em}record recursive call and current depth\\
\hspace*{1em}if propagation = baseline then\\
\hspace*{2em}(F', A', success) <- BCP\_BASELINE(F, A)\\
\hspace*{1em}else\\
\hspace*{2em}(state', success) <- BCP\_WATCHED(state)\\
\hspace*{1em}if success = false then return failure\\
\hspace*{1em}if all clauses are satisfied then return SAT with assignment\\
\hspace*{1em}l <- CHOOSE\_LITERAL(current formula, current assignment, heuristic)\\
\hspace*{1em}for decision literal d in [l, -l] do\\
\hspace*{2em}record one decision\\
\hspace*{2em}save current assignment/watch rollback marks\\
\hspace*{2em}apply d to the current branch state\\
\hspace*{2em}if branch immediately creates empty clause then\\
\hspace*{3em}record conflict, rollback branch state, and continue\\
\hspace*{2em}result <- DPLL(next branch state)\\
\hspace*{2em}if result is SAT then return result\\
\hspace*{2em}rollback branch state\\
\hspace*{2em}record one backtrack\\
\hspace*{1em}return failure
\end{quote}

At a high level, the baseline behavior is simple: try to force as much as
possible using propagation, then branch only when necessary. The branch order
and the efficiency of propagation are the two main levers that determine
performance, so these became the primary targets for improvement.

\subsection{Boolean Constraint Propagation}

BCP is the most frequently executed operation inside DPLL. In the baseline
solver, BCP repeatedly scans the current formula for unit clauses. When a unit
clause $(l)$ is found, the corresponding assignment is forced. Clauses
satisfied by $l$ are removed, occurrences of $\neg l$ are deleted from the
remaining clauses, and if any clause becomes empty, the branch is inconsistent.

\begin{quote}
\small
\ttfamily
\raggedright
BCP\_BASELINE(F, A)\\
\hspace*{1em}repeat\\
\hspace*{2em}U <- set of unit clauses in F\\
\hspace*{2em}if U is empty then return (F, A, true)\\
\hspace*{2em}for each unit literal u in U do\\
\hspace*{3em}if u contradicts an existing assignment then\\
\hspace*{4em}record conflict and return failure\\
\hspace*{3em}assign u in A\\
\hspace*{3em}record propagation\\
\hspace*{3em}remove all clauses containing u\\
\hspace*{3em}remove literal -u from all remaining clauses\\
\hspace*{3em}if some clause becomes empty then\\
\hspace*{4em}record conflict and return failure
\end{quote}

This approach is straightforward and correct, but it is inefficient because it
may scan large parts of the formula many times even when only a few clauses are
affected by the latest assignment.

\subsection{Watched Literals}

The first advanced improvement was watched literals. The idea is to avoid
rechecking every clause after each assignment. Instead, each clause watches one
literal if the clause is unit, or two literals if its length is at least two.
When a literal becomes false, only the clauses currently watching that literal
need to be revisited.

\begin{quote}
\small
\ttfamily
\raggedright
INIT\_WATCHES(F)\\
\hspace*{1em}for each clause C in F do\\
\hspace*{2em}if |C| = 1 then\\
\hspace*{3em}watch C[0] and enqueue C[0]\\
\hspace*{2em}else\\
\hspace*{3em}watch the first two literals of C
\end{quote}

\begin{quote}
\small
\ttfamily
\raggedright
BCP\_WATCHED(state)\\
\hspace*{1em}while unit queue is not empty do\\
\hspace*{2em}u <- pop a unit literal\\
\hspace*{2em}assign u if not already assigned\\
\hspace*{2em}for each clause watching -u do\\
\hspace*{3em}try to move that watch to another literal in the clause\\
\hspace*{3em}if a replacement literal exists and is not false then\\
\hspace*{4em}move the watch and continue\\
\hspace*{3em}else let w be the other watched literal\\
\hspace*{3em}if w is false then\\
\hspace*{4em}record conflict and return failure\\
\hspace*{3em}if w is unassigned then enqueue w as a new unit\\
\hspace*{1em}return success
\end{quote}

Watched literals do not change the search tree by themselves. The same branch
choices are made if the decision heuristic is unchanged. Their purpose is to
make BCP cheaper by revisiting only clauses that are actually affected by the
latest falsified literal.

\subsection{DLIS Decision Heuristic}

The second advanced improvement was the DLIS branching heuristic. Instead of
choosing the first unassigned literal encountered in the formula, DLIS counts
how many times each currently unassigned literal appears in clauses that are not
yet satisfied. The literal with the highest count is selected.

\begin{quote}
\small
\ttfamily
\raggedright
CHOOSE\_DLIS(F, A)\\
\hspace*{1em}initialize literal counts to zero\\
\hspace*{1em}for each clause C in F do\\
\hspace*{2em}if C is already satisfied under A then continue\\
\hspace*{2em}for each literal l in C do\\
\hspace*{3em}if l is still unassigned then count[l] <- count[l] + 1\\
\hspace*{1em}return literal with maximum count
\end{quote}

The intuition behind DLIS is that satisfying a frequently occurring literal is
likely to simplify many clauses at once. In practice this often decreases the
number of decisions, conflicts, and backtracks, especially on larger random
benchmarks. However, counting literal frequencies also adds overhead, so the
runtime effect must be measured experimentally rather than assumed.

\subsection{Data Structures and Design Choices}

The solver is written in Python, so the data structures were chosen to balance
clarity, correctness, and reasonable performance for the course project scale.

\subsubsection{Formula Representation}

Each literal is represented as a signed integer. For example, literal $x_7$ is
stored as \texttt{7}, and $\neg x_7$ is stored as \texttt{-7}. This matches the
DIMACS format directly and avoids an extra translation layer.

Each clause is stored as an immutable tuple of integers, and the whole formula
is stored as a tuple of clauses. The main reason for using tuples is that they
are simple, compact, and safe to share across recursive calls without accidental
in-place modification. In the baseline propagation path, simplification produces new tuples
when clauses are reduced, which keeps branch behavior easy to reason about.

\subsubsection{Assignment Representation}

The solver stores the current partial assignment in a dense array indexed by
variable id:
\begin{center}
\texttt{vals[var] = 1 for true, 0 for false, -1 for unassigned}
\end{center}

This representation is convenient because:
\begin{itemize}
    \item DIMACS variable ids are dense positive integers, so array indexing is
    natural,
    \item lookup by variable id is constant time without dictionary overhead,
    \item the same representation works for both baseline propagation and
    watched-literal propagation,
    \item a trail can record which variables were assigned in the current
    branch so they can be reset during rollback.
\end{itemize}

The helper function for literal evaluation interprets the sign of the literal
and returns true, false, or unassigned under the current \texttt{vals} array.

\subsubsection{Watched-Literal State}

The watched-literal implementation uses a dedicated state object containing:

\begin{itemize}
    \item \texttt{clauses}: the immutable tuple of original clauses,
    \item \texttt{vals}: the current partial assignment array,
    \item \texttt{watch\_pos}: for each clause, the indices of the two
    watched positions,
    \item \texttt{watch\_map}: a map from literal to the list of clause ids
    currently watching that literal,
    \item \texttt{root\_units} and branch queues: FIFO queues of literals waiting to be
    propagated.
\end{itemize}

These choices were made for specific reasons.

First, \texttt{watch\_pos} stores indices rather than literal values. This
makes it easy to move a watch from one position to another without rewriting the
clause itself.

Second, \texttt{watch\_map} is indexed by literal, so when a literal becomes
false we can immediately find the only clauses that might need attention. This
is the key efficiency advantage of watched literals.

Third, the unit queue is implemented using \texttt{deque}, which supports cheap
append and pop operations from the front. A FIFO queue is sufficient because
the correctness of unit propagation does not depend on a particular order.

Finally, the watched-literal path uses in-place mutation with explicit rollback
rather than copying the entire watched state at every recursive branch. The
solver keeps a variable assignment trail and a \texttt{watch\_moves} stack. When
a branch fails, it rolls back assignments to the saved trail mark and restores
watch positions to the saved watch-move mark. This avoids repeatedly copying
large watch lists while still keeping backtracking behavior deterministic.

\subsubsection{Instrumentation}

Solver statistics are stored in a small \texttt{SolverStats} record containing
decisions, propagations, conflicts, backtracks, recursive calls, maximum depth,
and runtime. This instrumentation was important for two reasons. First, it
allowed us to check that different solver variants produce the same SAT/UNSAT
classification. Second, it allowed us to separate search-tree effects
(decisions, conflicts, backtracks) from low-level implementation effects
(runtime).

\section{Experimental Methodology}

\subsection{Benchmark Set}

The final benchmark set contains 47 DIMACS CNF instances. It combines small
sanity-check inputs with larger SATLIB-style random instances.
In addition to this core comparison set, we also ran selected 200-variable
\texttt{uf200-*} and \texttt{uuf200-*} instances as larger stress tests.

\begin{center}
\small
\renewcommand{\arraystretch}{1.08}
\begin{tabular}{|p{2.2cm}|p{0.9cm}|p{1.0cm}|p{2.2cm}|p{6.0cm}|}
\hline
\textbf{Family} & \textbf{Count} & \textbf{Type} & \textbf{Size} & \textbf{Role in Evaluation} \\
\hline
\texttt{uf20-*} & 2 & SAT & 20 vars, 91 clauses & Small random SAT instances used to observe early backtracking behavior. \\
\hline
\texttt{uf50-*} & 20 & SAT & 50 vars, 218 clauses & Main larger SAT benchmark family. \\
\hline
\texttt{unsat*.cnf} & 4 & UNSAT & 2--13 vars, 4--34 clauses & Small handcrafted UNSAT checks used mainly for debugging and sanity testing. \\
\hline
\texttt{uuf50-*} & 20 & UNSAT & 50 vars, 218 clauses & Main larger UNSAT benchmark family. \\
\hline
\end{tabular}
\end{center}

This mixture is useful because it lets us see two different effects:

\begin{itemize}
    \item on tiny instances, implementation overhead can dominate and hide the
    benefit of more advanced heuristics;
    \item on larger instances, the search and propagation costs become large
    enough that the heuristics have room to help.
\end{itemize}

\subsection{Compared Variants}

For every benchmark, we executed the following three variants:

\begin{itemize}
    \item \textbf{baseline}: \texttt{--heuristic baseline --propagation baseline}
    \item \textbf{watched}: \texttt{--heuristic baseline --propagation watched}
    \item \textbf{dlis+watched}: \texttt{--heuristic dlis --propagation watched}
\end{itemize}

When no options are provided, \texttt{DPLL\_heuristic.py} defaults to
\texttt{dlis+watched}. The options above are used when we want to explicitly
separate the three solver variants.

This comparison isolates the contribution of each enhancement:

\begin{itemize}
    \item baseline vs. watched shows the effect of changing only the
    propagation data structure;
    \item watched vs. dlis+watched shows the effect of changing only the
    decision heuristic;
    \item baseline vs. dlis+watched shows the effect of both enhancements
    together.
\end{itemize}

\subsection{Metrics}

The main metrics collected were:
\begin{itemize}
    \item \textbf{decisions}: number of branch attempts,
    \item \textbf{conflicts}: number of contradictions detected,
    \item \textbf{backtracks}: number of failed recursive branches,
    \item \textbf{runtime}: total execution time for the instance.
\end{itemize}

We also tracked propagations, recursive calls, and maximum depth, but the four
metrics above are the most directly useful for comparing solver behavior.

Because the code is written in Python and many small instances solve in
microseconds, raw runtime on tiny inputs is noisy. For that reason, count-based
metrics such as decisions, conflicts, and backtracks are especially important:
they reveal whether a heuristic is actually changing the logical search, not
just the implementation overhead.

Runtime can vary slightly across runs due to system load and Python interpreter
overhead. Therefore, the reported runtime values should be interpreted as
representative measurements from one benchmark run, while the count-based
metrics such as decisions, conflicts, and backtracks are deterministic for a
fixed solver configuration.

\subsection{Correctness Checking}

Each benchmark belongs to a SAT or UNSAT directory, so the expected result is
known in advance. The benchmark runner checks every output against that expected
classification. All three variants solved all 47 inputs correctly, so the
reported performance differences are comparisons among correct solvers.

\subsection{Solver Interface}

The submitted solver entry point is \texttt{DPLL\_heuristic.py}. It provides a
command-line interface for solving DIMACS CNF files and for selecting the
heuristic configuration used in the experiments. A single-file grading-style
run is:

\begin{quote}
\small
\ttfamily
python3 DPLL\_heuristic.py benchmarks/sat/example-1.cnf
\end{quote}

This default command runs the final \texttt{dlis+watched} configuration.

This interface accepts one or more input files as command-line arguments. In
the default grading-style mode, it prints only the SAT/UNSAT result and a
satisfying assignment when the instance is satisfiable. A typical output is:

\begin{quote}
\small
\ttfamily
\raggedright
RESULT:SAT\\
ASSIGNMENT:1=0 2=0 3=1 4=0 5=1 6=1 ...
\end{quote}

The same executable is used for the advanced variants. For internal
experiments, the optional \texttt{--print-stats} flag reports the statistics
used in the tables:

\begin{quote}
\small
\ttfamily
python3 DPLL\_heuristic.py input.cnf --print-stats
\end{quote}

The benchmark runner calls the solver through its Python API to compare the
three variants under the same input set. This keeps the command-line interface
simple while avoiding subprocess overhead during benchmarking.

\section{Results}

\subsection{Overall Results}

\begin{center}
\small
\renewcommand{\arraystretch}{1.08}
\begin{tabular}{|l|l|r|r|r|r|}
\hline
\textbf{Group} & \textbf{Variant} & \textbf{Avg Decisions} & \textbf{Avg Conflicts} & \textbf{Avg Backtracks} & \textbf{Avg Runtime(ms)} \\
\hline
Overall & baseline & 205.38 & 101.30 & 200.34 & 13.399 \\
Overall & watched & 205.38 & 101.30 & 200.34 & 5.962 \\
Overall & dlis+watched & 76.55 & 36.13 & 70.68 & 6.702 \\
\hline
SAT only & baseline & 97.70 & 44.96 & 87.39 & 5.766 \\
SAT only & watched & 97.70 & 44.96 & 87.39 & 2.771 \\
SAT only & dlis+watched & 40.96 & 15.04 & 28.96 & 3.760 \\
\hline
UNSAT only & baseline & 308.58 & 155.29 & 308.58 & 20.715 \\
UNSAT only & watched & 308.58 & 155.29 & 308.58 & 9.019 \\
UNSAT only & dlis+watched & 110.67 & 56.33 & 110.67 & 9.522 \\
\hline
\end{tabular}
\end{center}

These results show two immediate patterns.

First, \texttt{watched} leaves decisions, conflicts, and backtracks unchanged
but cuts runtime almost in half. This confirms that watched literals are acting
exactly as intended: they reduce the cost of propagation without changing the
logical search tree.

Second, \texttt{dlis+watched} dramatically decreases decisions, conflicts, and
backtracks. Relative to the baseline over the full benchmark set, it lowers:

\begin{itemize}
    \item decisions by \textbf{62.7\%},
    \item conflicts by \textbf{64.3\%},
    \item backtracks by \textbf{64.7\%},
    \item runtime by \textbf{50.0\%}.
\end{itemize}

This is strong evidence that DLIS changes the search tree in a favorable way by
selecting branch literals that simplify the remaining formula more effectively.

\subsection{Effect of Watched Literals Alone}

The comparison between \texttt{baseline} and \texttt{watched} is especially
clean because the decision strategy is identical in both cases. Every change in
runtime can therefore be attributed to the propagation mechanism rather than to
a different sequence of branch choices.

Over all 47 benchmarks, watched literals reduce average runtime from 13.399 ms
to 5.962 ms, a \textbf{55.5\%} improvement. On SAT instances alone, the
runtime drops by 51.9\%, and on UNSAT instances it drops by 56.5\%.

The fact that decisions, conflicts, and backtracks remain exactly the same is
not a weakness; it is the expected signature of a successful propagation data
structure optimization. Watched literals do not decide \emph{where} to search.
They only reduce the cost of maintaining consistency after each assignment.

\subsection{Effect of DLIS on Top of Watched Literals}

The comparison between \texttt{watched} and \texttt{dlis+watched} isolates the
decision heuristic. Here the propagation mechanism is the same, so differences
are caused by the branch literal choice.

On the SAT subset, DLIS reduces decisions from 97.70 to 40.96 on average, a
\textbf{58.1\%} reduction. Conflicts drop by 66.5\%, and backtracks drop by
66.9\%. Runtime is still 34.8\% lower than the baseline, although it is higher
than watched literals alone because recomputing DLIS counts adds overhead.

On the UNSAT subset, DLIS reduces decisions from 308.58 to 110.67 on average,
a \textbf{64.1\%} reduction. Conflicts drop by 63.7\%, backtracks by 64.1\%,
and runtime is 54.0\% lower than the baseline. As with the SAT subset, the
runtime difference between \texttt{watched} and \texttt{dlis+watched} is small
because DLIS saves search work but spends extra time counting literals.

This behavior is consistent with the intended role of DLIS. By choosing a
literal that appears frequently in unsatisfied clauses, the solver often
reaches contradictions faster in UNSAT cases and satisfying assignments with
fewer failed branches in SAT cases.

\subsection{Effect of Problem Size}

To analyze how runtime varies with problem size, it is useful to separate the
benchmarks into small sanity-check cases and larger random instances.

\begin{center}
\small
\renewcommand{\arraystretch}{1.08}
\begin{tabular}{|l|r|r|r|r|r|}
\hline
\textbf{Group} & \textbf{Count} & \textbf{Variant} & \textbf{Avg Decisions} & \textbf{Avg Backtracks} & \textbf{Avg Runtime(ms)} \\
\hline
Small SAT & 3 & baseline & 16.67 & 10.00 & 0.322 \\
Small SAT & 3 & watched & 16.67 & 10.00 & 0.382 \\
Small SAT & 3 & dlis+watched & 8.67 & 3.67 & 0.746 \\
\hline
Large SAT & 20 & baseline & 109.85 & 99.00 & 6.705 \\
Large SAT & 20 & watched & 109.85 & 99.00 & 3.768 \\
Large SAT & 20 & dlis+watched & 45.80 & 32.75 & 4.837 \\
\hline
Small UNSAT & 4 & baseline & 9.50 & 9.50 & 0.056 \\
Small UNSAT & 4 & watched & 9.50 & 9.50 & 0.104 \\
Small UNSAT & 4 & dlis+watched & 9.50 & 9.50 & 0.169 \\
\hline
Large UNSAT & 20 & baseline & 368.40 & 368.40 & 25.014 \\
Large UNSAT & 20 & watched & 368.40 & 368.40 & 10.948 \\
Large UNSAT & 20 & dlis+watched & 130.90 & 130.90 & 12.206 \\
\hline
\end{tabular}
\end{center}

This table reveals an important trend.

For the \textbf{small} SAT and UNSAT groups, the advanced variants can be
slower even when they reduce the search tree. For example, small SAT runtime
increases from 0.322 ms in the baseline to 0.746 ms in \texttt{dlis+watched}
even though decisions drop from 16.67 to 8.67. The reason is simple: on tiny
inputs the overhead of maintaining watch structures and recomputing DLIS counts
can dominate the total runtime.

For the \textbf{large} benchmark groups, the pattern reverses. On large SAT
instances, watched literals reduce runtime by \textbf{43.8\%}, and
\texttt{dlis+watched} reduces decisions by \textbf{58.3\%}. On large UNSAT
instances, watched literals reduce runtime by \textbf{56.2\%}, and
\texttt{dlis+watched} reduces decisions by \textbf{64.5\%} while still cutting
runtime by \textbf{51.2\%}.

This size-sensitive behavior is one of the most important conclusions of the
project. Advanced heuristics are not free: they introduce bookkeeping overhead.
Their benefit becomes clear only when the instance is large enough that
propagation and backtracking costs dominate the run time.

\subsection{Interpretation of SAT vs. UNSAT Behavior}

The solver behaves differently on SAT and UNSAT instances, which is typical for
DPLL-based search.

For SAT problems, once the solver reaches a satisfying assignment it can stop.
Therefore a good branching heuristic helps by guiding the search toward a model
quickly and avoiding dead-end branches. This is why DLIS significantly reduces
backtracks on the SAT subset.

For UNSAT problems, the solver must prove that \emph{every} branch fails. As a
result, the search trees are typically larger, and the baseline averages are
much higher: 308.58 decisions and 20.715 ms runtime on UNSAT, compared with
97.70 decisions and 5.766 ms on SAT. This also explains why watched literals
show particularly strong runtime improvement on UNSAT instances: when the search
tree is large, propagation happens more often, so a propagation optimization has
more opportunities to pay off.

\subsection{Limitations}

The project still has several limitations.

\begin{itemize}
    \item The solver is implemented in Python, so the absolute runtime is not
    competitive with industrial SAT solvers.
    \item The watched-literal implementation uses a trail and undo stack, but
    it is still a compact educational implementation rather than a fully
    optimized industrial memory layout.
    \item The solver does not include CDCL features such as learned clauses,
    non-chronological backtracking, VSIDS, clause database management, or
    restarts.
    \item The benchmark set is useful for educational comparison, but it is
    still much smaller and less diverse than standard solver competitions.
\end{itemize}

These limitations do not invalidate the experimental results, but they explain
why the solver should be interpreted as a carefully instrumented educational
implementation rather than a production-grade SAT engine.

\section{Conclusions}

This project successfully built a complete DPLL-based SAT solver and used it to
study the effect of two major enhancements: watched literals and DLIS. The
final solver classified all 47 benchmarks correctly and produced clear,
consistent trends.

The most significant conclusions are:

\begin{itemize}
    \item \textbf{BCP dominates solver work}, so improving propagation with
    watched literals yields a large and reliable runtime benefit.
    \item \textbf{Decision quality matters}. DLIS substantially reduces the
    number of decisions, conflicts, and backtracks, especially on larger
    instances.
    \item \textbf{Problem size matters}. On very small instances, heuristic
    overhead can outweigh the benefit; on larger instances, the improvements
    become much more visible.
    \item \textbf{The two enhancements help in different ways}. Watched
    literals primarily improve low-level propagation efficiency, while DLIS
    improves the shape of the search tree itself.
\end{itemize}

From an educational perspective, the project demonstrates an important general
lesson in SAT solving: even without full CDCL machinery, careful choices in
data structures and branching strategy can drastically improve the performance
of a basic DPLL solver. The final implementation therefore serves as both a
working solver and a concrete experimental demonstration of how classical SAT
solver ideas affect search behavior.

\end{document}
